\section*{Bab \ref{ch:g6:solutions-and-solubility} --- Larutan dan Kelarutan}

\begin{solution}{exo:g6:solutions-and-solubility:1}
\omterm{def:g6:solutions-and-solubility:solute}{Zat terlarut}: gula. \omterm{def:g6:solutions-and-solubility:solute}{Pelarut}: air. Ya, \omterm{def:g6:solutions-and-solubility:solute}{pelarutnya} air: air gula adalah
\omterm{def:g6:solutions-and-solubility:solute}{larutan berair}.
\end{solution}

\begin{solution}{exo:g6:solutions-and-solubility:2}
Larutan yang tidak dapat melarutkan lagi \omterm{def:g6:solutions-and-solubility:solute}{zat terlarutnya}. \omterm{def:g6:solutions-and-solubility:solute}{Zat terlarut}
yang ditambahkan ke dalamnya tetap tidak \omterm{def:g2:mixing-and-dissolving:dissolve}{larut} di dasar.
\end{solution}

\begin{solution}{exo:g6:solutions-and-solubility:3}
$250 + 20 = \qty{270}{g}$.
\end{solution}

\begin{solution}{exo:g6:solutions-and-solubility:4}
Etanol: \omterm{def:g6:solutions-and-solubility:miscible}{dapat bercampur}. Minyak goreng: \omterm{def:g6:solutions-and-solubility:miscible}{tidak dapat bercampur}. Cuka:
\omterm{def:g6:solutions-and-solubility:miscible}{dapat bercampur} (sebagian besar cuka berupa air).
\end{solution}

\begin{solution}{exo:g6:solutions-and-solubility:5}
Propanon (aseton). GHS02: sangat mudah menyala; GHS07: menyebabkan
iritasi mata, uapnya dapat menimbulkan kantuk atau pusing.
\end{solution}

\begin{solution}{exo:g6:solutions-and-solubility:6}
$360 \times \frac{500}{1000} = \qty{180}{g}$.
\end{solution}

\begin{solution}{exo:g6:solutions-and-solubility:7}
\qty{200}{g} air melarutkan paling banyak $360 \times \frac{200}{1000} =
\qty{72}{g}$. Tidak: $80 - 72 = \qty{8}{g}$ tertinggal di dasar.
\end{solution}

\begin{solution}{exo:g6:solutions-and-solubility:8}
Massa larutannya $225 + 25 = \qty{250}{g}$; $\frac{25}{250} =
\qty{10}{\percent}$.
\end{solution}

\begin{solution}{exo:g6:solutions-and-solubility:9}
Di gelas kimia pertama, tempat semua garam telah \omterm{def:g2:mixing-and-dissolving:dissolve}{larut} dan larutannya
belum jenuh.
\end{solution}

\begin{solution}{exo:g6:solutions-and-solubility:10}
Karbon dioksida dilarutkan di bawah tekanan dalam botol yang tertutup.
Membukanya menurunkan tekanan: air dapat menampung lebih sedikit gas,
dan gas yang berlebih keluar sebagai gelembung.
\end{solution}

\begin{solution}{exo:g6:solutions-and-solubility:11}
$2000 \div 0.04 = 50\,000$ kali lebih banyak.
\end{solution}

\begin{solution}{exo:g6:solutions-and-solubility:12}
Gas lebih sedikit \omterm{def:g2:mixing-and-dissolving:dissolve}{larut} dalam air hangat: kolam yang hangat mengandung
lebih sedikit oksigen terlarut, dan ikan-ikan kekurangan oksigen.
\end{solution}

\begin{solution}{pb:g6:solutions-and-solubility:1}
\textbf{1.} \omterm{def:g6:solutions-and-solubility:solute}{Zat terlarut}: garam. \omterm{def:g6:solutions-and-solubility:solute}{Pelarut}: air.

\textbf{2.} Air itu menampung garam sebanyak yang dapat ditampungnya:
tidak ada garam lagi yang dapat \omterm{def:g2:mixing-and-dissolving:dissolve}{larut} di dalamnya.

\textbf{3.} Garam yang tidak \omterm{def:g2:mixing-and-dissolving:dissolve}{larut} di dasar membuktikan bahwa air garam
itu jenuh: air itu menampung garam sebanyak mungkin.

\textbf{4.} $360 \times 20 = \qty{7200}{g} = \qty{7.2}{kg}$.

\textbf{5.} $20 + 7.2 = \qty{27.2}{kg}$.

\textbf{6.} $\frac{7.2}{27.2} \approx 0.26$: sekitar \qty{26}{\percent}.

\textbf{7.} Tidak: airnya dapat melarutkan \qty{7.2}{kg}. Air itu masih
dapat melarutkan $7.2 - 5 = \qty{2.2}{kg}$.

\textbf{8.} $20 - 5 = \qty{15}{kg}$ air (\qty{15}{L}).

\textbf{9.} $360 \times 15 = \qty{5400}{g} = \qty{5.4}{kg}$.

\textbf{10.} Garam itu keluar dari larutan sebagai kristal padat: garam
itu mengkristal dan turun ke dasar.

\textbf{11.} Kristal-kristal putih baru muncul dan menumpuk di dasar
bak.

\textbf{12.} $7.2 - 5.4 = \qty{1.8}{kg}$ garam telah mengkristal.
\end{solution}
